\section{Case Studies: Representative Candidates}
\label{sec:casestudies}

We present five representative candidates drawn from the three domain pairs
(Table~\ref{tab:candidate_summary}).
Candidates are described as \textit{structurally generated cross-domain connections
with variable mechanistic interpretability}, not as validated mechanistic hypotheses.
Each is classified as \textbf{positive control}, \textbf{weakly novel}, or
\textbf{artifact risk} based on the relation semantics audit (Run~014);
interpretation confidence is rated high / medium / low.

% Table 4: Candidate Validation Summary
% Source: paper_assets/candidate_validation_table.csv
\begin{table}[h]
\caption{%
  \textbf{Representative candidate validation summary (5 candidates, Runs~011--013).}
  Candidates are identified by subset, hop count, and alignment bridge used.
  Classification follows the Run~014 relation semantics audit.
  Interpretation confidence is determined by the weakest relation category
  in the path: DM-only paths $\to$ high; paths with CR or DN(key step) $\to$ low.
  Score = weighted rubric total (Eq.~\ref{eq:score}); all candidates scored
  $\geq$0.50 and passed the drift filter.
}
\label{tab:candidate_summary}
\small
\begin{tabularx}{\linewidth}{llcllX}
\toprule
ID & Path summary & Hops & Classification & Confidence & Key limitation \\
\midrule
A-1 & VHL$\to$HIF1A$\to$LDHA$\to$NADH$\to$r\_Ox. &
  5 & positive control & high &
  Known pathway (Warburg effect); not novel \\
\addlinespace
B-1 & PTGS1$\to$m\_AA$\to$r\_OxNat$\to$fg\_Catechol &
  3 & weakly novel & low &
  Direction inversion; AA$\to$catechol not established (CR edge) \\
\addlinespace
B-2 & PTGS2$\to$m\_AA$\to$r\_OxNat$\to$fg\_Catechol &
  3 & weakly novel & low &
  Same as B-1; COX-2 selectivity adds interest if direction resolved \\
\addlinespace
C-1 & TH$\to$m\_Dop$\to$r\_Methyl$\to$fg\_Piperidine &
  3 & artifact risk & low &
  Methylation of dopamine does not produce piperidine (CR edge) \\
\addlinespace
C-2 & TH$\to$m\_Dop$\to$MAOA$\to$m\_Ser$\to$r\_Deam. &
  4 & weakly novel & medium &
  Compartmentalisation: dopamine and serotonin in distinct neuron types \\
\bottomrule
\end{tabularx}

\smallskip
\noindent Sub.\ A: Cancer sig.\ / Metab.\ chem.\ (bridge: NADH).
Sub.\ B: Immunology / Nat.\ products (bridge: Arachidonic Acid).
Sub.\ C: Neurosci.\ / Neuro-pharma.\ (bridge: Dopamine).
Case studies for each candidate: Section~\ref{sec:casestudies}.
\end{table}


\subsection{Candidate A-1 --- VHL / HIF1A / LDHA / NADH Cascade (Positive Control)}
\label{cs:a1}

\textit{Classification: positive control.
Interpretation confidence: high.
Directionality risk: none. Chemistry validity risk: none.}

\smallskip
The pipeline's highest-quality Subset~A candidate connects the VHL tumour suppressor
to NADH oxidation chemistry via a 5-hop cross-domain path:
\[
  g\textunderscore\mathrm{VHL} \xrightarrow{\textit{encodes}}
  \mathrm{VHL} \xrightarrow{\textit{inhibits}}
  \mathrm{HIF1A} \xrightarrow{\textit{activates}}
  \mathrm{LDHA} \xrightarrow{\textit{req\_cofactor}}
  \mathrm{NADH} \xrightarrow{\textit{undergoes}}
  r\textunderscore\mathrm{Oxidation}
\]
The path traverses from the cancer-signaling domain (VHL, HIF1A, LDHA) to the
metabolic chemistry domain ($r$\_Oxidation) via NADH, the sole aligned bridge
in Subset~A.
Strong relations account for 60\% of the path; all filter guards pass with zero
drift flags.

This candidate reconstructs the VHL$\to$HIF-1$\alpha$$\to$Warburg axis in renal
cell carcinoma: VHL loss stabilises HIF-1$\alpha$, which upregulates LDHA,
leading to increased aerobic glycolysis and NADH oxidation demand.
The complete mechanism is well-characterised in the oncology literature.
The candidate's value is as a \textit{positive control}: the pipeline independently
recovers a known pathway from graph structure without explicitly encoding the Warburg
effect as a search target.
It is \textit{not} novel; classifying it as a novel hypothesis would be incorrect.

The structural redundancy of Subset~A's three filter-passing candidates
(all variations of the same VHL/HIF1A/LDHA cascade via different entry nodes)
reflects the low fan-out of the NADH bridge, and directly motivates Claim~3.

\subsection{Candidate B-1 --- PTGS1 / Arachidonic Acid / Catechol (Weakly Novel)}
\label{cs:b1}

\textit{Classification: weakly novel.
Interpretation confidence: low.
Directionality risk: high. Chemistry validity risk: high.}

\smallskip
Candidate B-1 connects COX-1 (PTGS1) to catechol functional-group chemistry
via a 3-hop cross-domain path through the aligned eicosanoid bridge:
\[
  \mathrm{PTGS1} \xrightarrow{\textit{catalyzes}}
  m\textunderscore\mathrm{AA} \xrightarrow{\textit{undergoes}}
  r\textunderscore\mathrm{OxidNat} \xrightarrow{\textit{produces}}
  fg\textunderscore\mathrm{Catechol}
\]
Strong-relation ratio: 0.667 (catalyzes + undergoes); the filter passes all four guards.

The cross-domain connection has partial biological support: catechol-containing plant
compounds (caffeic acid derivatives, catechins) are documented COX inhibitors.
However, two concurrent problems limit the interpretation.
First, \textit{directionality inversion}: the path implies COX-1 produces catechol
(forward direction), whereas the established pharmacological relationship is that
catechol compounds \textit{inhibit} COX-1.
Second, \textit{class-level chemistry}: the \texttt{produces} edge
($r$\_OxidationNat $\to$ $fg$\_Catechol) is a natural-products generalisation that
does not hold for arachidonic acid as the specific substrate;
the primary COX-1 product (prostaglandin G$_2$) is not a catechol compound.

Next validation steps: (1)~confirm whether any documented AA oxidation products
contain catechol structural motifs (ChEBI, LIPID MAPS); (2)~assess whether the
catechol-COX connection should be modelled as inhibition rather than production.

\subsection{Candidate B-2 --- PTGS2 / Arachidonic Acid / Catechol (Weakly Novel)}
\label{cs:b2}

\textit{Classification: weakly novel.
Interpretation confidence: low.
Directionality risk: high. Chemistry validity risk: high.}

\smallskip
Candidate B-2 is structurally identical to B-1 but substitutes the inducible
isoform COX-2 (PTGS2):
\[
  \mathrm{PTGS2} \xrightarrow{\textit{catalyzes}}
  m\textunderscore\mathrm{AA} \xrightarrow{\textit{undergoes}}
  r\textunderscore\mathrm{OxidNat} \xrightarrow{\textit{produces}}
  fg\textunderscore\mathrm{Catechol}
\]
The same directionality and chemistry validity limitations apply.
The pharmacological significance of COX-2 is higher: it is the primary target
for NSAIDs and the coxib class.
Several catechol-containing natural compounds (EGCG, quercetin, hydroxytyrosol)
show preferential COX-2 inhibition in published assay data.

The B-1/B-2 pair together constitute a testable selectivity hypothesis if the
catechol-AA connection is confirmed: whether catechol natural products show
differential activity at COX-1 vs.\ COX-2 as a function of structural reactivity
with AA-derived oxidation products.
Resolving the directionality inversion is prerequisite to treating either
candidate as actionable.

\subsection{Candidate C-1 --- TH / Dopamine / Methylation / Piperidine (Artifact Risk)}
\label{cs:c1}

\textit{Classification: artifact risk.
Interpretation confidence: low.
Directionality risk: none. Chemistry validity risk: high.}

\smallskip
Candidate C-1 connects tyrosine hydroxylase (TH) to piperidine-scaffold drug
chemistry via a 3-hop path:
\[
  \mathrm{TH} \xrightarrow{\textit{catalyzes}}
  m\textunderscore\mathrm{Dopamine} \xrightarrow{\textit{undergoes}}
  r\textunderscore\mathrm{Methylation} \xrightarrow{\textit{produces}}
  fg\textunderscore\mathrm{Piperidine}
\]
Strong-relation ratio: 0.667; all filter guards pass.
The pharmacological framing is plausible: piperidine-scaffold antipsychotics
(haloperidol, risperidone, pimozide) target dopamine receptors, and TH is the
rate-limiting enzyme in dopamine biosynthesis.

However, the \texttt{produces} edge
($r$\_Methylation $\to$ $fg$\_Piperidine) is the source of failure.
O-methylation of dopamine by COMT produces methoxydopamine, which is structurally
unrelated to piperidine.
The KG edge encodes a \textit{synthetic chemistry generalisation}
(methylation reactions are used in piperidine synthesis) rather than an in vivo
metabolic relationship.
This is an instance of the CR (chemically-risky) pattern in the relation
semantics audit: a reaction class $\to$ functional group class edge that does
not hold for the specific substrate in the path.

C-1 illustrates a critical filter precision gap: the filter correctly passes
\texttt{catalyzes}, \texttt{undergoes}, and \texttt{produces} as strong relations,
but cannot validate the chemical plausibility of the substrate--product pair.
C-1 is a KG artefact that passed the filter via a class-level generalisation.

\subsection{Candidate C-2 --- TH / Dopamine / MAOA / Serotonin / Deamination (Weakly Novel)}
\label{cs:c2}

\textit{Classification: weakly novel.
Interpretation confidence: medium.
Directionality risk: moderate. Chemistry validity risk: none. Context dependence risk: high.}

\smallskip
Candidate C-2 connects TH to serotonin deamination via MAO-A in a 4-hop path:
\[
  \mathrm{TH} \xrightarrow{\textit{catalyzes}}
  m\textunderscore\mathrm{Dop} \xrightarrow{\textit{is\_substrate\_of}}
  \mathrm{MAOA} \xrightarrow{\textit{catalyzes}}
  m\textunderscore\mathrm{Ser} \xrightarrow{\textit{undergoes}}
  r\textunderscore\mathrm{Deamination}
\]
Strong-relation ratio: 0.50; all filter guards pass.
Each individual edge represents established biochemistry: TH produces dopamine,
dopamine is an MAO-A substrate, MAO-A processes serotonin, serotonin is deaminated
to 5-HIAA.

The cross-pathway coupling hypothesis---that TH activity modulates serotonin
deamination via dopamine-serotonin competition at MAO-A---is mechanistically
derivable from MAO-A's known dual substrate specificity.
MAO inhibitors (MAOIs, used as antidepressants) affect both dopaminergic and
serotonergic metabolism, consistent with this coupling.

The principal limitation is \textit{cellular compartmentalisation}:
dopaminergic and serotonergic neurons are distinct cell populations.
The KG merges dopaminergic (TH, dopamine) and serotonergic (MAO-A, serotonin)
data without cell-type annotation, producing a path that is structurally valid
in the merged graph but requires co-localisation conditions not encoded in the KG.
The competition hypothesis would most plausibly apply in neurons co-expressing
both pathways (rare), in pharmacological MAOI conditions, or in in vitro assays.
C-2 is classified as weakly novel: the cross-pathway coupling is plausible but
the compartmentalisation objection must be addressed before it is actionable.
